\section{Holy Fear and Persevering Hope}

The thought of judgment awakens fear because our acts matter eternally and because we have sinned. It awakens hope because the Lord is just, mercy has been offered, and grace can make us ready. The Catechism joins conversion, holy fear, commitment to justice, and the blessed hope of Christ's return (CCC 1041). We should allow these truths to remain together in prayer. Fear can shake us out of indifference; hope gives courage to obey the summons we have heard.

Holy fear directs us toward God. It acknowledges his majesty, grieves over offense against him, and takes his warnings seriously. It becomes fruitful in repentance and fidelity. A fear which repeatedly returns to an imagined catastrophe can instead absorb the attention needed for a duty. When this occurs, we can return to what Christ has given us to do: obey a known command, receive the sacraments, or care for someone whose need is present. Persistent anxiety may require patient pastoral help in making that return.

Francis de Sales places the two sentences of Matthew~25 beside one another. He dwells upon the terror of being sent away and then upon the blessed welcome by which Christ draws his servants to himself (I.14, considerations~5--6). The meditation does not end in the contemplation of terror. It ends in repentance and gratitude for the time and means of salvation. Christian fear finds its proper direction when it makes us more willing to receive the Savior, and Christian hope finds its proper seriousness when it makes us willing to change.

Hope must therefore be distinguished from an assumption that our present course will end well regardless of what we choose. We ask for perseverance because it is a gift, and we cooperate because the gift calls for a real response. Prayer, sacramental life, amendment, works of mercy, and faithful duty belong to that cooperation. A person who knows a grave sin needs confession and conversion, not an argument for postponement. A person who has honestly repented needs trust in mercy and renewed fidelity, not a private sentence of hopelessness.

Christ's delay is a mercy in which fidelity can grow. Another year may allow a wound to be repaired, a habit resisted, a duty embraced, and love made more generous. It may also reveal that a fault we hoped to leave behind still requires patient struggle. Repeated weakness should make us more dependent upon grace and more honest about the means we need. The Lord calls us to persevere in returning and obeying. We should not make discouragement another reason to abandon the good.

Nor is vigilance confined to moments when we feel close to eternity. Much of life proceeds through familiar work. We keep promises, earn an honest living, care for those entrusted to us, worship, and pray. The servant is blessed when the returning master finds him doing the task assigned (Mt 24:45--46). This gives spiritual weight to steadiness. Readiness can be quiet, practical, and unnoticed, because its measure is the will of the Lord rather than the intensity of the servant's feelings.

Through repeated fidelity we may begin to desire the holiness which would make the meeting with Christ joyful. Fear of admitting a fault can give way to willingness to confess it. Resentment of a duty can diminish as charity grows. This growth may be slow and interrupted by failures. It remains something to ask for in prayer, including when meditation on judgment chiefly awakens fear. The Lord can deepen our desire to be with him while teaching us to obey.

